\par\Needspace{6\baselineskip}\medskip
\SourceLeaf{224}\begin{center}\hypertarget{native-224-article161}{}\pdfbookmark[2]{161.}{article-161}161.\end{center}
Den \emph{ersten Fall} unterteilen wir in drei Unterfälle.

I. Fällt der Punkt $B$ mit $A$ oder mit dem Gegenpunkt zusammen, so gilt $\delta=0$ oder $=180^0$. Die Größen $\gamma$, $\epsilon'$, $\epsilon''$ und die Punkte $D'$, $D''$ bleiben unbestimmt, während $\gamma'$, $\gamma''$, $\epsilon$ und die Punkte $D$, $B^*$ bestimmt sind. Der Punkt $C$ fällt notwendig mit $A$ zusammen. Durch Überlegungen, die denen in Artikel 140 entsprechen, erhält man leicht die folgende Gleichung.
\[
0=n'\frac{\sin(z-\sigma)}{\sin z}\cdot\frac{R'\sin\delta'}{R''\sin\delta''}
\cdot\frac{\sin(A''D-\delta'')}{\sin(A'D-\delta'+\sigma)}-n''.
\]
Alles in den Artikeln 141, 142 Dargestellte lässt sich somit auf diesen Fall übertragen, sofern man $a=0$ setzt und $b$ unmittelbar aus Gleichung 12 in Artikel 140 bestimmt. Die Größen $z$, $r'$, $\dfrac{n'r'}{n}\mathpunct{\raisebox{.45ex}{,}}$
$\dfrac{n'r'}{n''}$ werden ebenso wie oben berechnet. Sobald $z$ und damit die Lage des Punktes $C'$ bekannt sind, lassen sich die Lage des Großkreises $CC'$, sein Schnittpunkt mit dem Großkreis $A''B''$, also der Punkt $C''$, und damit die Bögen $CC'$, $CC''$, $C'C''$ beziehungsweise $2f''$, $2f'$, $2f$ angeben. Daraus ergibt sich schließlich
\[
r=\frac{n'r'}{n}\cdot\frac{\sin2f}{\sin2f'}\mathpunct{\raisebox{.45ex}{,}}\qquad
r''=\frac{n'r'}{n''}\cdot\frac{\sin2f''}{\sin2f'}\mathpunct{\cdot}
\]

II. Fällt der Punkt $B''$ mit $A''$ oder mit dem Gegenpunkt zusammen, so lässt sich alles eben Dargestellte übertragen, sofern man alles, was sich auf den ersten Ort bezieht, mit dem vertauscht, was sich auf den dritten bezieht.

III. Etwas anders ist der Fall zu behandeln, in dem $B'$ entweder mit $A'$ oder mit dem Gegenpunkt zusammenfällt. Hier fällt der Punkt $C'$ mit $A'$ zusammen. Die Größen $\gamma'$, $\epsilon$, $\epsilon''$ und die Punkte $D$, $D''$, $B^*$ bleiben unbestimmt. Der Schnittpunkt des Großkreises $BB''$ mit der Ekliptik\footnote{Allgemeiner mit dem Großkreis $AA''$. Der Kürze halber betrachten wir hier nur den Fall, in dem die Ekliptik als Grundebene gewählt wird.} lässt sich jedoch angeben. Seine Länge sei $=l'+\pi$. Mit ähnlichen Überlegungen wie in Artikel 140 erhält man die Gleichung
\[
0=n\frac{R\sin\delta\sin(A''D'-\delta'')}{R''\sin\delta''\sin(AD'-\delta)}
+n'r'\frac{\sin\pi}{R''\sin(l''-l'-\pi)}+n''.
\]
Den Koeffizienten von $n$, der mit $a$ in Artikel 140 übereinstimmt, bezeichnen wir mit demselben \SourceLeaf{225}Zeichen $a$ und den Koeffizienten von $n'r'$ mit $\beta$. Dabei lässt sich $a$ selbst hier auch nach der Formel
\[
a=-\frac{R\sin(l'+\pi-l)}{R''\sin(l''-l'-\pi)}
\]
bestimmen. Somit gilt $0=an+\beta n'r'+n''$. Verbindet man diese Gleichung mit
\[
P=\frac{n''}{n}\mathpunct{\raisebox{.45ex}{,}}\qquad
Q=2\left(\frac{n+n''}{n'}-1\right)r'^3,
\]
so ergibt sich
\[
\frac{\beta(P+1)}{P+a}r'^4+r'^3+\tfrac12Q=0,
\]
woraus sich der Abstand $r'$ bestimmen lässt, außer wenn $\beta=0$ gilt. In diesem Fall ergäbe sich daraus lediglich $P=-a$. Auch wenn der Fall $\beta=0$ (der uns zum dritten, im nächsten Artikel zu behandelnden Fall führen würde) ausbleibt, ist $\beta$ stets eine sehr kleine Größe. Daher kann $P$ nur wenig von $-a$ abweichen. Offenbar wird damit die Bestimmung des Koeffizienten $\dfrac{\beta(P+1)}{P+a}$ sehr unsicher, sodass sich $r'$ nicht mit Genauigkeit bestimmen lässt.

Ferner erhält man $\dfrac{n'r'}{n}=-\dfrac{P+a}{\beta}\mathpunct{\raisebox{.45ex}{,}}$
$\dfrac{n'r'}{n''}=-\dfrac{P+a}{\beta P}\mathpunct{:}$ dann lassen sich ähnlich wie in Artikel 143 leicht die Gleichungen
\[
\begin{aligned}
r\sin\zeta&=\frac{n'r'}{n}\cdot\frac{\sin\gamma''}{\sin\epsilon'}\sin(l''-l')\\[.5em]
r''\sin\zeta''&=-\frac{n'r'}{n''}\cdot\frac{\sin\gamma}{\sin\epsilon'}\sin(l'-l)\\[.5em]
r\sin(\zeta-AD')&=r''P\frac{\sin\gamma''}{\sin\gamma}\sin(\zeta''-A''D'),
\end{aligned}
\]
herleiten. Verbindet man sie mit den Gleichungen VIII und IX aus Artikel 143, so lassen sich die Größen $r$, $\zeta$, $r''$, $\zeta''$ bestimmen. Die übrigen Rechenschritte entsprechen den oben beschriebenen.
